\section{분배 전략과 대화 분석}\label{sec:strategies}

\noindent\textcolor{red}{\textbf{예상!}} 아래는 아직 분석하기 전의 집필 자리이며 분배와 대화에서 확인할 항목을 적은 것으로, 모델에서 관찰된 전략이나 그 효과를 보고하는 내용은 아니다.

\subsection{분배 시점과 금액}

같은 양을 주더라도 언제 얼마씩 주는지에 따라 남는 여유가 달라지므로, 매 라운드 잔량을 맞추는 분배와 오답 뒤 벌금을 버틸 만큼 보충하는 분배가 나타나는지 확인할 예정이다.
이때 잦은 소액 이전과 적은 횟수의 큰 이전을 비교하고, 각 이전 앞뒤의 잔량과 예고된 벌금을 함께 읽어 분배 시점과 금액이 어떤 처지에서 정해지는지 살펴볼 예정이다.
그림 \ref{fig:arena-gini}에는 모델별 라운드 변화와 분배 전후의 지니 계수를 넣되 관측 시점과 비교할 자리는 집계 전에 확정한다.

\subsection{대화에 나타난 판단}

분배 이유를 확인할 때에는 응답과 대화에 적힌 이유를 실제 토큰 이동과 대조하여, 잔량을 맞추겠다는 말이나 필요한 만큼 주겠다는 말이 이후 분배와 맞는지 검토할 예정이다.
그 말만으로 모델의 내부 판단을 확정하지 않으며, 같은 이유를 말하고도 다른 금액을 주거나 말한 이전을 실행하지 않은 경우도 함께 확인할 예정이다.

\begin{figure*}[t]
\centering
\resizebox{\textwidth}{!}{%% Reserved layout only: no empirical or simulated curves are drawn.
\begin{tikzpicture}[x=1cm,y=1cm]
  \foreach \x/\heading in {0/{모델별 라운드 변화},8.7/{분배 전후 비교}} {
    \draw[black!35,dashed,rounded corners=2pt] (\x,0) rectangle ++(8.1,3.2);
    \node[sgfont=9,text=black!70] at (\x+4.05,2.65) {\heading};
    \node[sgfont=9,text=black!55,align=center] at (\x+4.05,1.45)
      {\textcolor{red}{예상!}\\[4pt]Gini 계수 그래프 자리\\[4pt]원시 기록 집계 후 삽입};
  }
\end{tikzpicture}
}
\caption{\textcolor{red}{\textbf{예상!}} 모델별 라운드 변화와 분배 전후의 지니 계수 그래프를 넣을 자리이다.}
\label{fig:arena-gini}
\end{figure*}
